% !TeX root = ../main.tex

\paragraph{Homework \#5 Nearly free electron model}

\begin{problem}[Compare Kronig-Penney model with the nearly free electron model]
  The Kronig-Penney model allows one to consider both the nearly-free-electron
  and tight-binding limits. Consider the (repulsive) Dirac-Kronig-Penney model.
  It gives the following implicit equation for the energy bands
  \[
    \cos(ka) = \cos(qa) + u\sin(qa)/qa,
  \]
  where $q = \sqrt{2mE}/\hbar$ and $u = \frac{mU_0a}{\hbar^2} > 0$. Assuming
  that $u \ll q$, find the solution of above equation near the Bragg points
  $k \approx \pm\frac\pi a$, $\pm\frac{2\pi}{a}$, $\ldots\,$. Compare your
  result with that of the nearly-free electron model. (Reminder: the $\approx$
  sign means that $k$ is close yet no equal to the wavenumber of the Bragg
  point.)
\end{problem}
\begin{solution}
  Near the Bragg points, assume
  \[
    ka = n\pi + \kappa a, \qq{and}
    qa = q_n a + \delta.
  \]
  with $|\kappa|$, $|\delta| \ll 1$, $q_n = n\pi/a$.
  Keep the 2nd order of $\kappa$ and $\delta$, the implict equation can be
  written as
  \[
    (-1)^n \ab[1 - \frac12(\kappa a)^2] = (-1)^n \ab(1 - \frac12\delta^2)
                                        + u\frac{(-1)^n\delta}{n\pi}.
  \]
  Then, we can have the solution of $\delta$
  \[
    \delta = u/n\pi \pm \sqrt{(u/n\pi)^2 + (\kappa a)^2}.
  \]
  Since $q = \sqrt{2mE}/\hbar$, then we have
  \[
    E = \hbar^2(q_n + \delta/a)^2/2m
      \approxeq E_n^0 + \hbar^2 q_n \delta/ma,
  \]
  where $E_n^0 = \frac{\hbar^2q_n^2}{2m}$.
  Substitute $q_n = n\pi/a$, and the solution of $\delta$ into the expression
  of $E$
  \[
    E = E_n^0 + \frac{n\pi\hbar^2\delta}{m}
  = E_n^0 + \frac{n\pi\hbar^2}{ma^2}
            \ab[\frac{u}{n\pi} \pm
            \sqrt{\ab(\frac{u}{n\pi})^2 + (\kappa a)^2}].
  \]
  Substitute $u = mU_0a/\hbar^2$ into the expression above
  \[
    E = E_n^0 + \frac{U_0}{a} \pm
        \sqrt{\ab(\frac{U_0}{a})^2 + \ab(\frac{n\pi\hbar^2\kappa}{ma})^2}
  \]
  with the band gap
  \[
    \Delta E_n = E_+ - E_-
  = 2\sqrt{\ab(\frac{U_0}{a})^2 + \ab(\frac{n\pi\hbar^2\kappa}{ma})^2}.
  \]
  When $\kappa a \to 0$, the energy band and the band gap become
  \[
    E = E_n^0 + U_0/a \pm U_0/a, \qq{and} \Delta E_n = 2U_0/a.
  \]
  The results perfectly match that of the nearly-free electron model.
\end{solution}

\begin{problem}[Nearly free electron model]
  Consider 2D electrons subject to a weak periodic potential
  \[
    U(x,y) = U_0 [\cos(2\pi x/a) + \cos(2\pi y/a)].
  \]
  \begin{enumext}
    \item Find the energy bands near the \emph{edges} of the first Brillouin
    zone. (Hint: the edges of the Brillouin zone are the Bragg planes where the
    eigenstates are doubly degenerate.) Plot the bands and isoenergetic
    surfaces.
    \item Repeat the analysis of part (a) for regions near the \emph{corners}
    of the first Brillouin zone, where there are four degenerate eigenstates.
  \end{enumext}
\end{problem}
\begin{solution}\leavevmode
  \begin{enumext}
    \item The reciprocal vectors are
    $\bm G_1 = \ab(\frac{2\pi}{a}, 0)$,
    $\bm G_2 = \ab(0, \frac{2\pi}{a})$.
    Apply Fourier transformation to the potential
    \[
      U(x,y) = \frac{U_0}{2}
      \ab[e^{\pm\iu\bm G_1 \cdot \bm r} + e^{\pm\iu\bm G_2 \cdot \bm r}]
      = \sum_G U_G \upe^{\iu\bm G \cdot \bm r}.
    \]
    So, the Fourier coefficient is $U_{\pm G_{1,2}} = U_0/2$.
    Near the zone edge (e.g., $k_x = \pi/a$), consider a wavevector
    \[
      \bm k = (\pi/a + q_x,  k_y)
    \]
    with $|q_x| \ll \pi/a$. The two nearly degenerate plane waves are
    $\ket|\bm k>$ and $\ket|\bm k - \bm G_1>$.
    Their free-electron energies are
    \begin{multline*}
      E^{(0)}(\bm k, \bm k - \bm G_1)
    = \frac{\hbar^2}{2m} \ab[\ab(\pm\frac{\pi}{a} + q_x)^2 + k_y^2]\\
    \approx E_0 \pm \alpha q_x
    = \frac{\hbar^2}{2m} \ab[\ab(\frac{\pi}{a})^2 + k_y^2]
      \pm \frac{\hbar^2 \pi}{m a} q_x.
    \end{multline*}
    The perturbation couples these states with matrix element
    $\braket<\bm k|U|\bm k - \bm G_1> = U_0/2$.
    In the subspace $\{\ket|\bm k>, \ket|\bm k - \bm G_1>\}$, the Hamiltonian is
    \[
      H = \begin{pNiceMatrix}[first-col, first-row]
                      & \ket|\bm k>      & \ket|\bm k - \bm G_1>\\
          \bra<\bm k| & E_0 + \alpha q_x & U_0/2 \\
\bra<\bm k - \bm G_1| & U_0/2 & E_0 - \alpha q_x
      \end{pNiceMatrix}.
    \]
    From the eigenfunction $\det(\mathcal H - E\identity) = 0$,
    one can obtain the eigenvalues
    \[
      E_{\pm} = E_0 + \frac{\hbar^2 q_x^2}{2m}
        \pm \sqrt{\ab(\frac{\hbar^2 \pi}{m a} q_x)^2 + \ab(\frac{U_0}{2})^2}.
    \]
    \item The corners ($M$-point) of the first Brillouin zone are
    \begin{align*}
      \bm k_1 & = \bm k_0 = (+\pi/a, +\pi/a), \\
      \bm k_2 & = \bm k_0 - \bm G_1 = (-\pi/a, +\pi/a), \\
      \bm k_3 & = \bm k_0 - \bm G_2 = (+\pi/a, -\pi/a), \\
      \bm k_4 & = \bm k_0 - \bm G_1 - \bm G_2 = (-\pi/a, -\pi/a),
    \end{align*}
    with the same free-electron energies
    \[
      E_0 = \frac{\hbar^2}{2m} \ab(k_x^2 + k_y^2) = \frac{\hbar^2 \pi^2}{m a^2}.
    \]
    In the basis
    $\{\ket|\bm k_1>, \ket|\bm k_2>, \ket|\bm k_3>, \ket|\bm k_4> \}$,
    states are coupled if they differ by $\pm \bm G_1$ or $\pm \bm G_2$.
    The non-zero off-diagonal matrix elements are
    \[
      \braket<\bm k_1|U|\bm k_2> = \braket<\bm k_1|U|\bm k_3> =
      \braket<\bm k_2|U|\bm k_4> = \braket<\bm k_3|U|\bm k_4> = \tfrac{U_0}2.
    \]
    The Hamiltonian is
    \[
      \mathcal H = E_0 \identity + \frac{U_0}{2}
      \begin{pNiceMatrix}[first-col, first-row]
      & \ket|\bm k_1> & \ket|\bm k_2>
      & \ket|\bm k_3> & \ket|\bm k_4> \\
        \bra<\bm k_1| & 0 & 1 & 1 & 0 \\
        \bra<\bm k_2| & 1 & 0 & 0 & 1 \\
        \bra<\bm k_3| & 1 & 0 & 0 & 1 \\
        \bra<\bm k_4| & 0 & 1 & 1 & 0
      \end{pNiceMatrix}.
    \]
    From the eigenfunction $\det(\mathcal H - E\identity) = 0$, one can obtain
    the eigenvalues $E = \{E_0 + U_0,\ E_0 - U_0,\ E_0,\ E_0\}$.
  \end{enumext}
\end{solution}